\label{sec:stand_der_forschung}

\subsection{Verfügbare Datenbasis und Korpora (Themenblock 1)}
\label{sec:datenbasis_korpora}

Einen guten Überblick über den aktuellen Stand und die verfügbaren Ressourcen für Leichte und Einfache Sprache in Deutschland geben \citet{madina2023easy}. Sie zeigen dabei vor allem auf, wie unterschiedlich die bestehenden Datensätze strukturiert sind und dass es an standardisierten Werkzeugen zur automatischen Bewertung der Texte fehlt.

Die Erforschung automatischer Textvereinfachung (Automatic Text Simplification, ATS) für das Deutsche litt lange Zeit unter einem deutlichen Mangel an Trainingsdaten. Erste Versuche, ein paralleles Korpus aufzubauen, gehen auf \citet{klaper2013klaper} mit dem \textit{KLAPER}-Korpus zurück, das jedoch mit ca. 250 parallelen Textpaaren (rund 7.000 Sätze und 70.000 Tokens) für moderne Deep-Learning-Ansätze deutlich zu klein war. 

Einen Meilenstein stellte die Arbeit von \citet{battisti2020corpus} dar, die ein umfangreiches, multimodales Korpus mit 6.217 Dokumenten (rund 211.000 Sätze, 5,5 Millionen Tokens) erstellten. Trotz dieser beachtlichen Größe zeigt der Datensatz eine klare Einschränkung: Während 5.461 Dokumente rein monolingual vorliegen, umfasst der parallele Teil lediglich 378 ausgerichtete Paare. Dieser kleine parallele Kern verdeutlicht, wie knapp ausgerichtete Trainingsdaten für überwachte Übersetzungsmodelle im Deutschen bisher waren.

Für den Bereich der "`Einfachen Sprache"' (Plain German) präsentierten \citet{stodden2023deplain} mit \textit{DEplain} ein professionell erstelltes paralleles Korpus. Dieses teilt sich auf in einen Nachrichtenbereich (DEplain-APA) mit ca. 500 Dokumentenpaaren (rund 13.000 manuell ausgerichteten Satzpaaren) und einen Webbereich mit ca. 150 manuell ausgerichteten Dokumenten (ca. 2.000 Satzpaaren). Um die Datenknappheit weiter zu reduzieren, stellt das Framework zusätzlich einen Web-Harvester zur automatischen Extraktion und Ausrichtung bereit, wodurch der Webbereich auf ca. 750 Dokumentenpaare (ca. 3.500 automatisch ausgerichtete Satzpaare) vergrößert werden kann. Die Autoren zeigten, dass Modelle, die auf Dokumentenebene trainiert werden, oft Probleme haben, die feinkörnigen semantischen Übergänge beizubehalten, weshalb präzise Alignments eine kritische Ressource für das Training von Seq2Seq-Modellen darstellen. An dieser Stelle ist jedoch hervorzuheben, dass eine direkte Vergleichbarkeit mit Systemen zur Übersetzung in Leichte Sprache erschwert ist: Obwohl das DEplain-Framework methodisch wertvolle Ansätze liefert, bezieht es sich durch den Fokus auf Einfache Sprache auf eine andere linguistische Ziel- und Datendomäne.

Einen weiteren Beitrag zur Generierung paralleler Ressourcen leisteten \citet{toborek2023new} mit der Vorstellung des \textit{New Aligned Simple German Corpus}. Ihr Datensatz besteht aus ca. 700 Dokumentenpaaren auf Dokumentenebene und über 10.000 Satzpaaren auf Satzebene, die teils manuell annotiert und teils automatisch berechnet wurden. Das Korpus kombiniert dabei das systematische Scraping paralleler Webinhalte mit automatischer Satz-Ausrichtung. Allerdings fasst die Arbeit Texte sowohl der Leichten als auch der Einfachen Sprache unter einem gemeinsamen Vereinfachungsbegriff zusammen, was eine feine Unterscheidung der spezifischen syntaktischen und semantischen Richtlinien von Leichter Sprache erschwert. Dieser Kompromiss verdeutlicht jedoch auch ein grundlegendes Problem der aktuellen Forschung: Ein Korpusaufbau, der sich ausschließlich auf reine Leichte Sprache beschränkt, ist Stand jetzt kaum in großem Umfang realisierbar, da es im deutschsprachigen Web schlichtweg an einer ausreichenden Anzahl frei zugänglicher, paralleler Quellen fehlt.

\subsection{Evaluierungsmetriken für die Vereinfachung (Themenblock 2)}
\label{sec:evaluierungsmetriken}

Die automatische Evaluierung von Textvereinfachungssystemen stellt eine große methodische Hürde dar. Traditionell stützt sich die automatisierte Lesbarkeitsbewertung im deutschen Sprachraum auf deterministische Formeln wie den \textit{Flesch Reading Ease} \citep{amstad1978verstaendlich}, die \textit{Wiener Sachtextformel} \citep{bamberger1984lesen} oder den \textit{LIX-Index} \citep{bjornsson1968readability} (deren mathematische Definitionen in Abschnitt~\ref{sec:linguistische_masse_lesbarkeit} dargelegt sind). Diese Metriken quantifizieren die Oberflächenkomplexität primär anhand von Silben-, Wort- und Satzlängenstatistiken. Während sie für allgemeine Lese- und Sachtexte eine robuste, recheneffiziente Orientierungsbasis bieten, sind sie konzeptionell nicht darauf ausgelegt, die spezifischen, normierten Regeln der Leichten Sprache (wie das Vermeiden von Passivkonstruktionen, strikte Hauptsatz-Parataxen oder die Bindestrich-Gliederung von Komposita) zu erfassen.

Klassische referenzbasierte N-Gramm-Metriken der maschinellen Übersetzung wie BLEU \citep{papineni2002bleu} und ROUGE \citep{lin2004rouge} messen die lexikalische Übereinstimmung mit einer Referenzübersetzung. Für die automatische Textvereinfachung erweisen sich diese jedoch oft als ungeeignet: Wie \citet{alva2021unsuitability} empirisch belegen, korreliert BLEU im Vereinfachungskontext nur schwach oder sogar negativ mit dem menschlichen Urteil über Einfachheit und Textqualität. Der Grund liegt darin, dass BLEU unzulässige 1:1-Kopien des komplexen Originals mit hohen Punktwerten belohnt, während legitime lexikalische Variationen und strukturelle Modifikationen (wie Satzteilungen und Umschreibungen) fälschlicherweise als Fehler abgestraft werden.

Als Alternative für die Textvereinfachung etablierte sich die Metrik \textit{SARI} (\textit{System Accessibility and Readability Index}) nach \citet{xu2016optimizing}. SARI vergleicht die Systemgenerierung $y$ simultan mit dem komplexen Ausgangstext $x$ und den Referenzübersetzungen $R$:
\begin{equation}
    \text{SARI} = \frac{F_{\text{add}} + F_{\text{keep}} + P_{\text{del}}}{3}
\end{equation}
Hierbei messen $F_{\text{add}}$ und $F_{\text{keep}}$ den F1-Score für erfolgreich eingefügte bzw. beibehaltene N-Gramme, während $P_{\text{del}}$ die Präzision für getilgte N-Gramme des komplexen Originals bewertet. Dadurch honoriert SARI explizit Vereinfachungsoperationen, die sich vom komplexen Ausgangstext wegbewegen.

Im Laufe der Entstehungszeit dieser Arbeit wurde mit \textit{DETECT} von \citet{korobeynikova2024detect} eine neuartige, erlernbare Evaluierungsmetrik speziell für die deutsche Textvereinfachung vorgestellt. DETECT adaptiert das englischsprachige LENS-Framework für den deutschsprachigen Raum und zielt darauf ab, die Qualität automatischer Vereinfachungen ganzheitlich über drei Dimensionen abzubilden: Verständlichkeit (\textit{Simplicity / Ease}), semantische Inhaltsbewahrung (\textit{Meaning Preservation}) sowie grammatikalische Korrektheit und Sprachfluss (\textit{Fluency / Textual Clarity}). Um den Mangel an manuell annotierten Trainingsdaten zu überwinden, wird das Regressionsmodell vollständig auf synthetisch durch Large Language Models generierten Qualitätsurteilen trainiert. In empirischen Validierungen auf einem eigens erhobenen menschlichen Testdatensatz wiesen die Autoren nach, dass DETECT deutlich stärkere Korrelationen mit menschlichen Qualitätsbewertungen erzielt als traditionelle Maße wie BLEU, SARI oder BERTScore.

Gleichzeitig verdeutlicht dieser Ansatz zentrale verbleibende Forschungslücken: Da der zugrundeliegende Evaluierungsdatensatz (\textit{SimpEvalDE}) primär auf Korpora der Einfachen Sprache (DEplain-APA und LHA-APA) aufbaut, sind spezifische, normierte Richtlinien der Leichten Sprache — wie das strikte Vermeiden von Passiv- und Genitivkonstruktionen oder die typografische Gliederung von Komposita — nicht explizit als Kriterien modelliert. Zudem sind rein referenzbasierte Satz-Metriken nur bedingt geeignet, um die Komplexität längerer Dokumente ganzheitlich und referenzfrei zu bewerten oder als kontinuierliches Steuerungssignal während des Modelltrainings eingesetzt zu werden.

Obwohl das Thema Leichte Sprache in der deutschsprachigen NLP-Forschung über viele Jahre hinweg nur wenig Beachtung gefunden hat, zeigt das jüngste Erscheinen solcher Arbeiten, dass das Forschungsfeld derzeit stark an Dynamik gewinnt, was die wissenschaftliche und gesellschaftliche Relevanz dieses Themenbereichs zusätzlich unterstreicht.

\subsection{Maschinelle Übersetzung und Optimierungsverfahren (Themenblock 3)}
\label{sec:maschinelle_textvereinfachung_alignment}
\label{sec:maschinelle_uebersetzung_modellierung}
\label{sec:dpo_textgenerierung}

Die maschinelle Übersetzung von Alltagssprache (AS) in Leichte Sprache (LS) lässt sich formal als monolinguale maschinelle Übersetzung (\textit{Monolingual Machine Translation}) begreifen. Im Gegensatz zur klassischen bilingualen Übersetzung erfordert dieser Transformationsprozess tiefgreifende asymmetrische Strukturänderungen: Komplexe Satzgefüge müssen aufgespalten (\textit{Sentence Splitting}), Passiv- in Aktivkonstruktionen überführt, irrelevante Textpassagen gekürzt (\textit{Compression}) und Fachbegriffe explizierend umschrieben werden.

Erste systematische Benchmarks neuronaler Übersetzungsmodelle für deutsche Leichte Sprache von \citet{saeuberli2020benchmarking} zeigten, dass Sequence-to-Sequence-Modelle grundlegende Satzkürzungen und Wortsubstitutionen erlernen können, bei komplexen syntaktischen Umstellungen jedoch stark durch die geringe Menge paralleler Trainingsdaten limitiert werden.

Bisherige Ansätze operieren jedoch fast ausnahmslos auf isolierter Satzebene und stützen sich auf manuell annotierte Präferenzdaten. Eine zentrale Forschungslücke besteht darin, wie sich generative Übersetzungsmodelle auf Dokumenten- und Artikelebene trainieren und durch eine automatisierte Reward-Steuerung über erlernte Komplexitätsmetriken optimieren lassen.

